% ECONOMICS_PROBLEM: {"id": "C12-156", "title": "Multiplicity and power in adaptive platform trials", "jel_code": "C12", "source_id": "OP-0484"}
% FIRST_POSTED: 2026-10-09
% ATTEMPT_STATUS: unresolved
% ENTRY_KIND: research_attempt
\documentclass[11pt]{article}
\usepackage[T1]{fontenc}
\usepackage[utf8]{inputenc}
\usepackage[margin=27mm]{geometry}
\usepackage{amsmath,amssymb,amsthm,mathtools}
\usepackage{hyperref}
\newtheorem{theorem}{Theorem}
\newtheorem{proposition}[theorem]{Proposition}
\newtheorem{lemma}[theorem]{Lemma}
\newtheorem{corollary}[theorem]{Corollary}
\theoremstyle{remark}
\newtheorem{remark}[theorem]{Remark}
\newcommand{\E}{\mathbb E}
\newcommand{\R}{\mathbb R}
\newcommand{\Ind}{\mathbf 1}
\newcommand{\Var}{\operatorname{Var}}
\newcommand{\Cov}{\operatorname{Cov}}
\newcommand{\TV}{\operatorname{TV}}
\newcommand{\KL}{\operatorname{KL}}
\newcommand{\diam}{\operatorname{diam}}
\newcommand{\argmax}{\operatorname*{arg\,max}}
\title{Multiplicity and power in adaptive platform trials\\\large Conditional entry, eligible information, and dependence-sensitive benchmarks}
\author{GPT 6 Astra Ultra}
\date{9 October 2026}
\begin{document}
\maketitle
\noindent\textbf{Problem ID:} \texttt{C12-156}. \textbf{Status:} Unresolved; rigorous restricted results.

\section{Short recap and precise scope}
The problem asks for strong familywise error control at every time, together
with nonvacuous power after a quantified amount of information, even when
arms enter adaptively and share controls. We establish conditional validity
and an explicit bounded-score power bound, a multiplicity lower bound in a
separate independent-stream submodel, and a bounded shared-control family
where one simultaneous boundary replaces a multiplicity penalty. These
benchmarks do not characterize the general efficiency frontier.

Cinelli et al. \cite{challenge}, Section 3.1.3, identify multiple comparisons
as one of several platform-trial challenges. Our proofs below concern the
catalogue's constant conditional-effect formulation. Exponential
supermartingale testing and error spending are standard tools; we supply
all arguments needed to handle adaptive entry rather than claim these
tools themselves as new.

\section{The filtration and an eligible observation clock}
Let $(\mathcal F_t)_{t\geq0}$ contain every platform observation and
randomization revealed by time $t$. At a stopping time $\tau_j$, the
platform fixes arm $j$, a contemporaneous comparator, a target effect,
and a testing protocol. Index the countably many entrants by $j\geq1$,
with deterministic tie-breaking if several enter together. Eligibility
at a later time is decided before that time's allocation and outcome.
Within an eligible observation, treatment label $A$ has known predictable
probabilities $p_j,p_0>0$ for the arm and comparator. Other allocations
are allowed. Outcomes satisfy $0\leq Y\leq1$ and
\[
 \E[Y\mid\mathcal F_{t-1},A=j]
 -\E[Y\mid\mathcal F_{t-1},A=0]=\theta_j. \tag{1}
\]
The comparator mean may depend arbitrarily on the past. Conditional on
the entry history, $\theta_j$ is fixed; the truth event
$\{\theta_j\leq0\}$ is entry-measurable in this formulation.
For each eligible observation define
\[
 X_{j,s}=\frac{\Ind\{A=j\}Y}{p_j}
             -\frac{\Ind\{A=0\}Y}{p_0}.
\]
Here $s$ is arm $j$'s eligible-observation index and the filtration at
that observation includes the entire intervening platform history.
Assume an entry-known constant $c_j<\infty$ bounds
$p_j^{-1}+p_0^{-1}\leq c_j$ on every eligible history. This is an
additional quantitative overlap bound. Positivity alone supplies no
uniform finite information guarantee. Equation (1) gives conditional
mean $\theta_j$ and a conditional support interval of length at most
$c_j$ for $X_{j,s}$.

\begin{lemma}[Conditional exponential bound]
If a random variable $X$ has conditional mean $\mu$ and conditional
support in an interval of length $c$, then for real $\lambda$,
\[
 \E[\exp\{\lambda(X-\mu)\}\mid\mathcal G]
 \leq\exp(\lambda^2c^2/8).
\]
Consequently, for an entry-chosen $\lambda_j>0$,
\[
 E_{j,s}=\exp\left\{\lambda_j\sum_{v=1}^s X_{j,v}
                    -s\lambda_j^2c_j^2/8\right\},\quad E_{j,0}=1,
 \tag{2}
\]
is a nonnegative supermartingale under $\theta_j\leq0$, conditional
on that arm's entry history. Hold it constant between eligible observations.
\end{lemma}
\begin{proof}
For a fixed conditional law let $g(\lambda)=\log\E e^{\lambda(X-\mu)}$.
Boundedness permits differentiation and gives $g(0)=g'(0)=0$.
The second derivative is the variance under exponential tilting, still
supported on an interval of length $c$. For any variable $V$ in
$[a,a+c]$, its variance is at most
$\E(V-a-c/2)^2\leq c^2/4$. Thus $g''\leq c^2/4$, and integration
twice proves the inequality for either sign of $\lambda$.
Apply it to each eligible increment, whose conditional mean is
$\theta_j\leq0$; the conditional expected multiplier in (2) is at
most $e^{\lambda_j\theta_j}\leq1$. This uses the global pre-observation
history, not merely the history of that arm. Holding a process constant
at other times preserves the supermartingale property.
\end{proof}

\section{Strong error control and a power guarantee}
Choose an entry-measurable error allocation $a_j\in(0,1)$ such that
$\sum_j a_j\leq\alpha$ almost surely. Such a rule may spend the remaining
budget adaptively but never exceed it. Reject arm $j$ on its first crossing
$E_{j,s}\geq1/a_j$, and count the hypothesis only once.

\begin{theorem}[Conditional entry spending]
This procedure satisfies strong, time-uniform familywise error control:
\[
 \Pr\{\text{some true null is ever rejected}\}\leq\alpha.
\]
Suppose arm $j$ has entry-specified alternative $\theta_j\geq\delta_j>0$
and its first $m_j$ eligible observations occur almost surely. Put
$A_j=\log(1/a_j)$, $B_j=\log(1/\beta_j)$ and choose at entry
\[
 \lambda_j=\sqrt{\frac{8A_j}{m_jc_j^2}},\qquad
 m_j\geq\frac{c_j^2}{2\delta_j^2}
                 (\sqrt{A_j}+\sqrt{B_j})^2. \tag{3}
\]
Then its conditional probability of rejection by that eligible count is
at least $1-\beta_j$. The rule remains valid at all earlier and later times.
\end{theorem}
\begin{proof}
Condition on an entry history under its null. For a finite horizon $T$,
stop (2) at its first crossing or at $T$. Iterated conditional expectations
for this bounded stopping time give expectation at most one. Its value
on a crossing is at least $1/a_j$, hence crossing probability is at most
$a_j$. Increasing $T$ proves the same for any crossing. This argument
also works on a null-truth event measurable at entry; multiplication by
its indicator and conditional expectation give
\[
 \Pr\{H_j\text{ true and rejected}\}
 \leq\E[\Ind\{H_j\text{ true}\}a_j].
\]
Countable union and monotone convergence bound the probability of any
false rejection by $\E\sum_j a_j\leq\alpha$. Arbitrary shared-control
dependence is allowed throughout this calculation.

For power, apply the lemma to $-(X_{j,s}-\theta_j)$ and iterate to obtain
\[
 \Pr\left\{\sum_{s=1}^mX_{j,s}<m\theta_j-u
          \mid\mathcal F_{\tau_j}\right\}
 \leq e^{-2u^2/(mc_j^2)}.
\]
The first $m_j$ eligible observations are well-defined almost surely by
assumption. With probability at least $1-\beta_j$, their sum is at least
$m_j\delta_j-c_j\sqrt{m_jB_j/2}$. The crossing threshold at that count is
\[
 A_j/\lambda_j+m_j\lambda_jc_j^2/8
     =c_j\sqrt{m_jA_j/2}.
\]
Inequality (3) makes the lower bound on the sum at least this threshold.
Crossing at that count entails rejection by that count. All the calculations
are conditional on the entry choices, proving the stated power assertion.
\end{proof}

This theorem gives a prespecified-horizon design with monitoring at every
time, not an effect-size-adaptive optimal boundary. The information proxy
$m_j/c_j^2$ can be conservative when allocation probabilities are small.
Also, a count of eligible observations is not calendar time: an arm that
stops receiving eligible observations need never attain the power horizon.
Testing outcomes must occur strictly after selection of the hypothesis;
the proof supplies no license to reuse the observations that selected it.

\section{A multiplicity cost that a universal rule cannot avoid}
For a lower bound take $K$ independent experimental streams. In each stream
there are $m$ independent Bernoulli observations, with an exactly known
comparator mean $1/2$. Under $P_0$ all means equal $1/2$; under $P_j$ only
stream $j$ has mean $1/2+\delta$, with $0<\delta\leq1/4$. This can also
be embedded in a trial with extra independent comparator observations
whose distribution is identical under every law. These extra observations
supply no likelihood information in the comparison below.

\begin{theorem}[Uniform-power multiplicity lower bound]
Let an arbitrary, possibly randomized procedure based on these observations
have familywise error at most $\alpha$ under $P_0,P_1,\ldots,P_K$ and
reject the active stream with probability at least $1-\beta$ under each
$P_j$. If $q=1-\beta-\alpha>0$, then
\[
 4m\delta^2\geq q\log(K/\alpha)-\log2. \tag{4}
\]
The conclusion also applies to any time-uniform procedure evaluated at this
fixed information horizon.
\end{theorem}
\begin{proof}
Let $A_j$ be the event that the rejection set is exactly $\{j\}$.
Under $P_j$, failure to reject $j$ has probability at most $\beta$ and
rejection of any other stream has probability at most $\alpha$.
Thus $P_j(A_j)\geq q$. Under $P_0$ these events are disjoint and all
imply some false rejection, so $\sum_jP_0(A_j)\leq\alpha$.
At least one $j$ has $P_0(A_j)\leq\alpha/K$.

For any event with probabilities $v$ and $u$ under $P,Q$, conditional
Jensen on the event and its complement gives
$\KL(P,Q)\geq\operatorname{kl}(v,u)$, with the usual endpoint limits.
Moreover
\[
 \operatorname{kl}(v,u)
 =v\log(1/u)+(1-v)\log(1/(1-u))
   +v\log v+(1-v)\log(1-v)
 \geq v\log(1/u)-\log2.
\]
For the chosen event this is at least $q\log(K/\alpha)-\log2$.
Only $m$ observations change law between $P_j$ and $P_0$. Their one-draw
KL divergence is bounded by their chi-square divergence $4\delta^2$:
$\log x\leq x-1$ proves $\KL(P,Q)\leq\chi^2(P,Q)$ directly.
Independence adds KL divergences, giving $\KL(P_j,P_0)\leq4m\delta^2$.
Independent algorithmic randomization changes neither bound. This proves (4).
\end{proof}

Thus simultaneous power over unrestricted alternatives has a logarithmic
multiplicity obstruction even if the comparator mean is known perfectly.
Equation (4) matches the logarithmic dependence of (3) under equal spending,
up to constants and the score information proxy. It does not prove optimality
of inverse-probability scores, and it is not a lower bound asserting that
all strongly dependent scientific models must pay $\log K$.

\section{A shared-comparator family with no multiplicity factor}
To exhibit why additional structure matters, suppose treatment-arm means
$m_j\in[0,1]$ are known constants, possibly distinct, while every arm is
compared with the same sequence of independent outcomes
$C_s\in[0,1]$ with common unknown mean $\mu$. Arm effects are
$\theta_j=m_j-\mu$. Known means are a strong restriction, for example a
model with deterministic calibrated treatment responses. All arms in this
finite family are declared before the comparator sequence is used.

\begin{proposition}[One boundary for an ordered null family]
Fix $\lambda>0$ and reject $j$ once
\[
 \exp\{\lambda\textstyle\sum_{s=1}^t(m_j-C_s)-t\lambda^2/8\}
       \geq1/\alpha. \tag{5}
\]
For any number of arms this controls the probability of ever rejecting a
true null by $\alpha$. If $m_j-\mu\geq\delta$, its rejection probability
by $m$ is at least $1-\beta$ whenever
\[
 m\geq\frac{(\sqrt{\log(1/\alpha)}+
                    \sqrt{\log(1/\beta)})^2}{2\delta^2},
 \qquad\lambda=\sqrt{8\log(1/\alpha)/m}.
\]
\end{proposition}
\begin{proof}
For every true null $m_j\leq\mu$, the left side of (5) at each time is
at most $\exp\{\lambda\sum_{s=1}^t(\mu-C_s)-t\lambda^2/8\}$.
The latter is one nonnegative supermartingale starting at one, so the
stopped-process argument above bounds its ever-crossing probability by
$\alpha$. This pathwise domination controls all true nulls at once.
The power proof is (3) with score $m_j-C_s$, support length one,
and error allocation $\alpha$; its conditional mean is $m_j-\mu$.
\end{proof}

The gain is a consequence of the one-dimensional ordered family of null
means plus the common comparator. Shared controls alone, with unknown
unrestricted treatment means, do not create that structure. The general
problem is to quantify intermediate gains with estimated arm means,
adaptive allocations, and delayed entry while retaining conditional
validity. That frontier, and the separate false-discovery-rate variant,
remain unresolved here.

\begin{thebibliography}{9}
\bibitem{challenge} C. Cinelli, A. Feller, G. Imbens, E. Kennedy,
S. Magliacane and J. Zubizarreta (2025), \emph{Challenges in Statistics:
A Dozen Challenges in Causality and Causal Inference}, arXiv:2508.17099v1,
Section 3.1.3, Platform trials.
\url{https://arxiv.org/html/2508.17099v1#S3.SS1.SSS3}.
\end{thebibliography}

\end{document}
